\section{Detection Playbooks}

Detection playbooks translate knowledge of adversary behavior into concrete, repeatable logic that can run continuously 
in Splunk or Microsoft Sentinel. Unlike hunting playbooks, which are exploratory and analyst-driven, detection playbooks 
are designed for automation: they define what to look for, how to detect it, how to investigate it, and how to tune it 
over time. As noted in the source document, “these playbooks form the backbone of a mature detection program and serve 
as a long-term reference for building and maintaining high-fidelity alerts.” 

Each detection playbook follows a consistent structure: a description of the behavior and why it matters, required telemetry, 
SPL and KQL detection logic, Linux investigation pivots, and tuning considerations. This ensures detections are both 
actionable and adaptable across environments.

% --------------------------------------------------------------------
% PLAYBOOK 1: Encoded PowerShell Commands
% --------------------------------------------------------------------

\subsection*{Detection Playbook: Encoded PowerShell Commands}

\textbf{Behavior and Rationale}

Adversaries frequently abuse PowerShell to execute malicious code, download payloads, or perform reconnaissance. 
To evade simple string-based detections, they often use the \texttt{-EncodedCommand} parameter, which allows Base64-encoded 
commands. As the source document states, “this behavior is highly suspicious in most environments and is a strong candidate 
for a high-priority detection.” 

\textbf{Required Telemetry}

Sysmon Event ID 1 or Windows Event ID 4688 with full command-line logging.  
Splunk: \texttt{sysmon}, \texttt{wineventlog}  
Sentinel: \texttt{Sysmon}, \texttt{SecurityEvent}

\textbf{SPL Detection Logic}

\begin{lstlisting}
index=sysmon EventCode=1 Image="*powershell.exe"
| search CommandLine="*EncodedCommand*"
| stats count by Computer, User, ParentImage, Image, CommandLine
| where count >= 1
\end{lstlisting}

\textbf{KQL Detection Logic}

\begin{lstlisting}
Sysmon
| where EventID == 1 and Image endswith "powershell.exe"
| where CommandLine contains "EncodedCommand"
| summarize Count=count() by Computer, User, ParentImage, Image, CommandLine
| where Count >= 1
\end{lstlisting}

\textbf{Linux Pivot}

\begin{lstlisting}[language=bash]
grep -i "powershell" /var/log/sysmon/sysmon.log | grep -i "EncodedCommand"
\end{lstlisting}

\textbf{Tuning Considerations}

Some administrative scripts legitimately use encoded PowerShell. Tune by excluding:

\begin{itemize}
    \item known management accounts
    \item automation servers
    \item approved scripts
\end{itemize}

All other occurrences should remain high-priority.

% --------------------------------------------------------------------
% PLAYBOOK 2: Excessive Failed Logons
% --------------------------------------------------------------------

\subsection*{Detection Playbook: Excessive Failed Logons}

\textbf{Behavior and Rationale}

Repeated failed logons indicate brute-force attacks, password spraying, or credential guessing. The document notes that 
“patterns of repeated failures against specific accounts or from specific hosts are highly suspicious.” 

\textbf{Required Telemetry}

Windows: Event ID 4625  
Linux: \texttt{/var/log/auth.log}

\textbf{SPL Detection Logic}

\begin{lstlisting}
index=wineventlog EventCode=4625
| stats count by AccountName, WorkstationName
| where count > 10
\end{lstlisting}

\textbf{KQL Detection Logic}

\begin{lstlisting}
SecurityEvent
| where EventID == 4625
| summarize Count=count() by Account, Workstation
| where Count > 10
\end{lstlisting}

\textbf{Linux Pivot}

\begin{lstlisting}[language=bash]
grep "Failed password" /var/log/auth.log | awk '{print $9}' | sort | uniq -c
\end{lstlisting}

\textbf{Tuning Considerations}

False positives may come from:

\begin{itemize}
    \item misconfigured services
    \item outdated credentials in scripts
    \item users repeatedly entering incorrect passwords
\end{itemize}

Tune by excluding noisy systems and service accounts.

% --------------------------------------------------------------------
% PLAYBOOK 3: Rare Parent-Child Process Relationships
% --------------------------------------------------------------------

\subsection*{Detection Playbook: Rare Parent-Child Process Relationships}

\textbf{Behavior and Rationale}

Legitimate processes follow predictable parent-child relationships. Rare combinations may indicate process injection, 
masquerading, or malicious execution. The document explains that “detecting rare parent-child combinations is a powerful 
behavioral technique.” 

\textbf{Required Telemetry}

Sysmon Event ID 1 with \texttt{ParentImage} and \texttt{Image} fields.

\textbf{SPL Detection Logic}

\begin{lstlisting}
index=sysmon EventCode=1
| stats count by ParentImage, Image
| where count < 3
\end{lstlisting}

\textbf{KQL Detection Logic}

\begin{lstlisting}
Sysmon
| where EventID == 1
| summarize Count=count() by ParentImage, Image
| where Count < 3
\end{lstlisting}

\textbf{Linux Pivot}

\begin{lstlisting}[language=bash]
ps aux --forest
\end{lstlisting}

\textbf{Tuning Considerations}

Rare does not always mean malicious. Tune by:

\begin{itemize}
    \item building a baseline of known-good rare pairs
    \item excluding expected administrative tasks
    \item focusing on scripting engines and interpreters
\end{itemize}

% --------------------------------------------------------------------
% PLAYBOOK 4: Rare Outbound Network Connections
% --------------------------------------------------------------------

\subsection*{Detection Playbook: Rare Outbound Network Connections}

\textbf{Behavior and Rationale}

Compromised hosts often communicate with external command-and-control servers. These connections are typically low-frequency 
and target unusual ports. The document states that “detecting rare outbound connections helps identify such activity.” 

\textbf{Required Telemetry}

Sysmon Event ID 3, firewall logs, DNS logs.

\textbf{SPL Detection Logic}

\begin{lstlisting}
index=sysmon EventCode=3
| stats count by DestinationIp, DestinationPort
| where count < 5
\end{lstlisting}

\textbf{KQL Detection Logic}

\begin{lstlisting}
Sysmon
| where EventID == 3
| summarize Count=count() by DestinationIp, DestinationPort
| where Count < 5
\end{lstlisting}

\textbf{Linux Pivot}

\begin{lstlisting}[language=bash]
ss -tulnp
\end{lstlisting}

\textbf{Tuning Considerations}

Tune by excluding:

\begin{itemize}
    \item known vendor support IPs
    \item internal infrastructure
    \item expected cloud service endpoints
\end{itemize}

Focus on unknown external IPs and unusual ports.

% --------------------------------------------------------------------
% PLAYBOOK 5: Persistence via Scheduled Tasks and Cron Jobs
% --------------------------------------------------------------------

\subsection*{Detection Playbook: Persistence via Scheduled Tasks and Cron Jobs}

\textbf{Behavior and Rationale}

Scheduled tasks (Windows) and cron jobs (Linux) are common persistence mechanisms. Detecting their creation or modification 
is critical for identifying long-term footholds.

\textbf{Required Telemetry}

Windows: Event ID 4698  
Linux: cron configurations in \texttt{/etc/crontab}, \texttt{/etc/cron.d/}, user cron files.

\textbf{SPL Detection Logic}

\begin{lstlisting}
index=wineventlog EventCode=4698
| table _time, Computer, TaskName, Author, Command
\end{lstlisting}

\textbf{KQL Detection Logic}

\begin{lstlisting}
SecurityEvent
| where EventID == 4698
| project TimeGenerated, Computer, TaskName, SubjectUserName, Command
\end{lstlisting}

\textbf{Linux Pivot}

\begin{lstlisting}[language=bash]
cat /etc/crontab
ls -al /etc/cron.d/
crontab -l
\end{lstlisting}

\textbf{Tuning Considerations}

Tune by excluding:

\begin{itemize}
    \item backup jobs
    \item monitoring scripts
    \item known administrative tasks
\end{itemize}

Focus on tasks invoking scripting engines or binaries from unusual paths.

% --------------------------------------------------------------------
% SUMMARY
% --------------------------------------------------------------------

\subsection*{Summary}

The detection playbooks in this section illustrate how to move from abstract knowledge of adversary behavior to concrete, 
automated logic in Splunk and Sentinel. As the document states, “each playbook combines behavioral focus, clear telemetry 
requirements, SPL and KQL implementations, Linux investigation pivots, and practical tuning guidance.” 

Over time, these playbooks can be expanded, refined, and adapted to new threats, forming a living detection library that 
supports long-term growth as a CTI analyst, threat hunter, and detection engineer.
